\chapter[Índices]{\emj{🔍}~Índices}

\section[¿Qué es un índice?]{\emj{🧠}~¿Qué es un índice?}

Un índice es una estructura de datos auxiliar que mantiene una copia ordenada de uno o más atributos de una tabla, con punteros a las filas originales en el heap (la tabla principal).
La analogía del índice de un libro es precisa: en lugar de leer todo el libro para encontrar un tema (Seq Scan), vas directamente a la página referenciada (Index Scan).

Sin índice, el motor debe leer \emph{cada fila} de la tabla para encontrar las que cumplen el filtro — esto es $O(n)$.
Con un índice B-tree, la búsqueda es $O(\log n)$ para localizar la primera fila, y luego secuencial para las siguientes.
En una tabla de 10 millones de filas, la diferencia entre Seq Scan e Index Scan puede ser de 8 segundos vs 0.3 milisegundos.

\begin{teoriabox}{Estructura interna de un índice B-tree}
Un B-tree (árbol B) es un árbol balanceado de búsqueda con estas propiedades:
\begin{itemize}
  \item \textbf{Nodos internos:} contienen claves y punteros a nodos hijo. Guían la búsqueda de $O(\log n)$.
  \item \textbf{Nodos hoja:} contienen las claves ordenadas y el puntero al heap (tid: page + offset). Son el nivel que finalmente referencia los datos.
  \item \textbf{Balanceado:} todas las hojas están a la misma profundidad. Un árbol con 1 millón de entradas tiene solo ~20 niveles.
  \item \textbf{Soporte de rangos:} las hojas están enlazadas entre sí, permitiendo escaneos de rango eficientes (BETWEEN, >, <).
\end{itemize}
\textbf{Costo real de mantener un índice:}
cada INSERT, UPDATE o DELETE debe actualizar todas las estructuras de índice de la tabla.
En una tabla con 5 índices, cada escritura genera 5 operaciones adicionales de árbol.
En tablas con escritura masiva, esto puede ser el cuello de botella.
\end{teoriabox}

\section[Tipos de índices en PostgreSQL]{\emj{📋}~Tipos de índices en PostgreSQL}

PostgreSQL ofrece varios tipos de índices porque no existe una estructura óptima para todos los patrones de consulta.
Cada tipo está optimizado para un conjunto específico de operadores y tipos de datos.
Elegir el tipo incorrecto puede hacer que el índice sea ignorado por el planificador.

\begin{itemize}
  \item \textbf{B-tree (default):} árbol balanceado. Soporta: \texttt{=, <, >, <=, >=, BETWEEN, LIKE 'prefix\%', IS NULL, IN}.
    El 90\% de los casos. Si no sabes qué tipo usar, usa B-tree.

  \item \textbf{Hash:} tabla hash. Solo soporta \texttt{=}. Ligeramente más rápido que B-tree para igualdad exacta, pero no soporta rangos y no se replica en versiones antiguas de PG. Poco usado.

  \item \textbf{GIN} (Generalized Inverted Index): índice invertido — mapea elementos a las filas que los contienen.
    Ideal para: arrays (\texttt{@>, \&\&}), JSONB (\texttt{@>, ?}), full-text search (\texttt{@@}).
    Lento en escritura (reconstruye la estructura), muy rápido en consultas.

  \item \textbf{GiST} (Generalized Search Tree): framework extensible para tipos de datos no estándar.
    Usado para: geometría/geo (PostGIS), rangos (\texttt{tsrange, daterange}), full-text.
    Más versátil que GIN, algo más lento en búsqueda.

  \item \textbf{BRIN} (Block Range Index): almacena rangos mínimo-máximo por bloque físico.
    Muy compacto (tamaño ínfimo comparado con B-tree). Eficiente solo cuando los datos están físicamente ordenados por la columna indexada (series de tiempo insertadas en orden, logs).
\end{itemize}

\section[Cuándo crear (y no crear) un índice]{\emj{💡}~Cuándo crear (y no crear) un índice}

La \textbf{selectividad} es el factor más importante: un índice es útil cuando filtra una fracción pequeña de las filas.
Si una condición retorna el 50\% de las filas, un Seq Scan puede ser más eficiente que un Index Scan porque lee el heap secuencialmente (más amigable con el cache de disco).
El planificador de PostgreSQL conoce estadísticas de distribución de valores y decide automáticamente cuándo usar el índice.

\begin{tipbox}{Reglas prácticas para índices}
\textbf{Crea índice en:}
\begin{itemize}
  \item Columnas en \texttt{WHERE} con alta selectividad (pocos resultados de muchas filas)
  \item Columnas de FK — PostgreSQL \emph{no} las indexa automáticamente; MySQL sí
  \item Columnas en \texttt{JOIN ON} y \texttt{ORDER BY} en queries frecuentes con tablas grandes
  \item Columnas de búsqueda full-text (índice GIN con \texttt{to\_tsvector})
\end{itemize}
\textbf{No crees índice en:}
\begin{itemize}
  \item Tablas con menos de 1{,}000 filas — el planificador preferirá Seq Scan de todas formas
  \item Columnas con cardinalidad muy baja: booleans, \texttt{status} con 3 valores posibles
  \item Columnas que se actualizan muy frecuentemente — el costo de mantenimiento supera la ganancia de lectura
  \item Columnas raramente usadas en \texttt{WHERE} — el índice ocupa espacio y ralentiza escrituras sin beneficio
\end{itemize}
\end{tipbox}

\begin{lstlisting}[caption={Crear índices B-tree, parcial y compuesto}]
-- Índice simple (B-tree por defecto)
CREATE INDEX idx_empleados_depto ON empleados(id_depto);

-- Índice compuesto
CREATE INDEX idx_pedidos_cliente_fecha ON pedidos(id_cliente, fecha DESC);

-- Índice parcial: solo indexa subset de filas
CREATE INDEX idx_pedidos_pendientes ON pedidos(fecha)
WHERE estado = 'pendiente';

-- Índice UNIQUE
CREATE UNIQUE INDEX idx_clientes_email ON clientes(email);

-- Crear índice SIN bloquear escrituras (producción)
CREATE INDEX CONCURRENTLY idx_ventas_fecha ON ventas(fecha);
\end{lstlisting}

\begin{lstlisting}[caption={EXPLAIN ANALYZE --- verificar uso de índice}]
-- Sin índice: Seq Scan
EXPLAIN ANALYZE
SELECT * FROM empleados WHERE salario > 80000;
-- Seq Scan on empleados (cost=0..520 rows=15)
-- Execution time: 8.4 ms

-- Crear índice
CREATE INDEX idx_empleados_salario ON empleados(salario);

-- Con índice: Index Scan
EXPLAIN ANALYZE
SELECT * FROM empleados WHERE salario > 80000;
-- Index Scan using idx_empleados_salario
-- Execution time: 0.3 ms

-- GIN para búsqueda en JSONB
CREATE INDEX idx_datos_gin ON tabla USING GIN (datos_jsonb);
SELECT * FROM tabla WHERE datos_jsonb @> '{"activo": true}';
\end{lstlisting}

\begin{lstlisting}[caption={Índices y FK en PostgreSQL}]
-- PostgreSQL NO crea índice automático en FK
-- Sin este índice, DELETE en tabla padre hace Seq Scan en hijos
CREATE INDEX idx_empleados_id_depto ON empleados(id_depto);
CREATE INDEX idx_pedidos_id_cliente ON pedidos(id_cliente);

-- Ver índices existentes
SELECT indexname, indexdef
FROM   pg_indexes
WHERE  tablename = 'empleados';
\end{lstlisting}

\begin{entrevistabox}
\begin{enumerate}
  \item ¿Cuándo usarías un índice GIN en lugar de B-tree?
  \item ¿Por qué los índices ralentizan las escrituras?
  \item ¿Qué es un índice parcial y cuándo tiene sentido usarlo?
\end{enumerate}
\end{entrevistabox}
